\chapter{Homogeneous objectives without compactness}\label{comp:H}\label{H:chapter}
\renewcommand{\ModuleTag}{7}
Coercive homogeneous objectives provide a broad global convergence family. This chapter goes beyond that compact case: a different potential gives noncompact value alternatives, recurrent-gradient obstructions and conditional iterate conclusions, which remain distinct from finite line-search escape and from whole-iterate convergence.


% SOURCE_BEGIN H.span.1
\begingroup

\setcounter{section}{0}
\setcounter{theorem}{0}
\setcounter{lemma}{0}
\setcounter{corollary}{0}
\setcounter{proposition}{0}
\setcounter{remark}{0}
\setcounter{equation}{0}
% SOURCE_END H.span.1


% SOURCE_BEGIN H.span.2
\section{Scope and main
conclusions}\label{H:scope-and-main-conclusions}

This chapter continues the investigation of ordinary, full-memory BFGS
under the algorithmic formulation in Xie--W\"achter~\cite{XW}. The accepted
steps satisfy fixed Armijo and \textbf{weak} Wolfe inequalities.
Accepted iterates are differentiable; the inverse matrices are positive
definite; no damping, restarts, smoothing, spectral safeguards, or
particular line-search implementation is imposed.

The principal family is

\[
f(x)=\max_{1\le i\le m}a_i^\top x,\qquad m<\infty.
\]

The objective-value theorem below holds more generally for every finite
convex, positively homogeneous function on the whole space. ``Finite''
in this broader statement means real-valued everywhere, not a finite
number of pieces.

\textbf{Proof status.} The value and iterate statements below have reviewed Lean proofs under their stated assumptions. Their formal hypotheses and dependencies are indexed with the source claims. Finite numerical tests and publication-priority assertions are not premises of those proofs. The reusable conditional lifting implications retain any projected-run limit explicitly stated in their hypotheses.

% SOURCE_END H.span.2


% SOURCE_BEGIN H.span.3
\subsection{Theorem 7.A: an unconditional objective-value
dichotomy}\label{H:theorem-a-an-unconditional-objective-value-dichotomy}

Let \(f:\mathbb R^n\to\mathbb R\) be finite, convex, and positively
homogeneous of degree one. Every infinite admissible full-BFGS run
satisfies

\[
\boxed{f(x_k)\longrightarrow0\quad\text{or}\quad f(x_k)\longrightarrow-\infty.}
\]

There is \textbf{no assumption} that the iterates, displacements,
accepted step lengths, or matrices are bounded. There is also no
assumption of a finite gradient alphabet for this theorem.

Consequently, if \(f\) is bounded below, then \(f(x_k)\to0=\min f\). If
\(f\) is coercive, then \(x_k\to0\), its unique minimizer. If an
infinite run ever attains \(f(x_k)\le0\), its values must tend to
\(-\infty\); the finite-value branch has \(f(x_k)>0\) at every
iteration.

% SOURCE_END H.span.3


% SOURCE_BEGIN H.span.4
\subsection{Theorem 7.B: a many-piece extension of the common-tilt
geometry}\label{H:theorem-b-a-many-piece-extension-of-the-common-tilt-geometry}

For a finite max-linear objective, let \(G\) be the set of gradients
occurring infinitely often. If the objective values have a finite limit,
then

\[
\boxed{0\in\operatorname{aff}G.}
\]

This no longer needs a bounded-displacement hypothesis. In particular,

\[
0\notin\operatorname{aff}\{a_1,\ldots,a_m\}
\quad\Longrightarrow\quad
f(x_k)\longrightarrow-\infty
\]

for every infinite admissible run. A useful arbitrary-many-piece family
is

\[
\boxed{f(u,t)=t+\max_{1\le i\le m}b_i^\top u.}
\]

All its gradients have last component one, so every infinite admissible
run drives the values to \(-\infty\). This is an infinite-run result,
not a proof of finite termination for every many-piece objective.

% SOURCE_END H.span.4


% SOURCE_BEGIN H.span.5
\subsection{Theorem 7.C: complete iterate dichotomy in the planar
unbounded-below
case}\label{H:theorem-c-complete-iterate-dichotomy-in-the-planar-unbounded-below-case}

For a finite max-linear objective on \(\mathbb R^2\) with
\(\inf f=-\infty\), every infinite admissible run satisfies exactly one
of

\[
\boxed{x_k\longrightarrow0,\quad f(x_k)\downarrow0,}
\]

or

\[
\boxed{f(x_k)\longrightarrow-\infty,\quad\|x_k\|\longrightarrow\infty.}
\]

In the first alternative the limit is nonstationary. The previously
constructed exact three-piece example attains that alternative. In the
second, convergence of the \emph{directions} \(x_k/\|x_k\|\) is not
asserted.

More generally, for any finite max-linear objective, if the span of the
recurrent gradients has dimension at most two, every finite-value-limit
run has convergent iterates. Section 7.7 proves this broader statement,
including the noncoercive bounded-below case.

% SOURCE_END H.span.5


% SOURCE_BEGIN H.span.6
\subsection{What this does not yet
establish}\label{H:what-this-does-not-yet-establish}

The objective-value classification is now unconditional in every
dimension. The remaining unresolved iterate question in this
investigation is the zero-value branch for noncompact objectives when
the recurrent gradient span has dimension at least three. We have
neither proved that all such iterate sequences converge nor constructed
a nonconvergent zero-value-limit trajectory.

This distinction is essential: the results close the previous
unbounded-displacement gap for \textbf{function values}, but are not
being presented as a full all-dimensional classification of the
\textbf{iterates}.


\setcounter{section}{1}
\setcounter{theorem}{0}
\setcounter{lemma}{0}
\setcounter{corollary}{0}
\setcounter{proposition}{0}
\setcounter{remark}{0}
\setcounter{equation}{0}
% SOURCE_END H.span.6


% SOURCE_BEGIN H.span.7
\section{Notation and BFGS
identities}\label{H:notation-and-bfgs-identities}

Write

\[
\begin{aligned}
g_k&=\nabla f(x_k),&p_k&=-H_kg_k,&s_k&=\alpha_kp_k,\\
x_{k+1}&=x_k+s_k,&y_k&=g_{k+1}-g_k,&B_k&=H_k^{-1},\\
q_k&=g_k^\top H_kg_k>0,&d_k&=-g_k^\top s_k=\alpha_kq_k>0,\\
t_k&=y_k^\top s_k,&\Delta_k&=f(x_k)-f(x_{k+1}).
\end{aligned}
\]

The line-search conditions are

\[
f(x_{k+1})\le f(x_k)+c_1g_k^\top s_k,
\qquad
g_{k+1}^\top s_k\ge c_2g_k^\top s_k,
\qquad0<c_1<c_2<1.
\]

In particular, with \(c=1-c_2>0\),

\[
t_k\ge c\,d_k>0,
\qquad
\beta_k:=\frac{t_k}{d_k}
=-\frac{y_k^\top H_kg_k}{q_k}\ge c.
\]

The full inverse-BFGS update is

\[
H_{k+1}=
\left(I-\frac{s_ky_k^\top}{t_k}\right)
H_k
\left(I-\frac{y_ks_k^\top}{t_k}\right)
+\frac{s_ks_k^\top}{t_k}.
\]

Its Hessian form, after substituting \(s_k=-\alpha_kH_kg_k\), is

\[
\boxed{B_{k+1}=B_k-\frac{g_kg_k^\top}{q_k}
+\frac{y_ky_k^\top}{t_k}.}
\tag{7.2.1}
\]

A zero gradient corresponds to a stationary termination rather than an
infinite strict-descent run. Throughout, ``infinite admissible''
therefore includes nonzero accepted gradients and well-defined
positive-curvature updates.

\Needspace{19\baselineskip}


\setcounter{section}{4}
\setcounter{theorem}{0}
\setcounter{lemma}{0}
\setcounter{corollary}{0}
\setcounter{proposition}{0}
\setcounter{remark}{0}
\setcounter{equation}{0}
% SOURCE_END H.span.7


% SOURCE_BEGIN H.span.8
\section{Homogeneous potential and proof of Theorem
7.A}\label{H:homogeneous-potential-and-proof-of-theorem-a}

% SOURCE_END H.span.8


% SOURCE_BEGIN H.span.9
\subsection{Bounded subgradients and Euler's
identity}\label{H:bounded-subgradients-and-eulers-identity}

For finite convex positively homogeneous \(f\), every subgradient
\(g\in\partial f(x)\) satisfies

\[
g^\top x=f(x),\qquad g^\top z\le f(z)\quad\text{for every }z.
\]

The first equality follows by applying the subgradient inequality to
\(0\) and \(2x\); the second follows by substituting that equality back
into the inequality. In particular, all subgradients belong to
\(\partial f(0)\). In positive dimension, the unit sphere is nonempty and \(f\) is finite and continuous there, so
\[
\|g\|\le\max_{\|z\|=1}f(z)<\infty.
\]
In dimension zero the space is a singleton and every gradient is zero; boundedness and the later bounded-sublevel or coercive convergence conclusions hold directly. All uses of a sphere minimum or maximum below concern positive dimension. This leaves the value theorem in arbitrary finite dimension; an infinite admissible run with nonzero gradients cannot occur in dimension zero.

Thus the bounded-gradient assumption in Theorem~\ref{M:average} is automatic, even for
an infinite-piece sublinear function.

% SOURCE_END H.span.9


% SOURCE_BEGIN H.span.10
\subsection{The exact energy
identity}\label{H:the-exact-energy-identity}

Set \(f_k=f(x_k)\) and

\[
h_k=y_k^\top x_{k+1}=d_k-\Delta_k,\qquad
\ell_k=-y_k^\top x_k.
\]

Convexity and homogeneity give \(h_k,\ell_k\ge0\) and
\(t_k=h_k+\ell_k\). Armijo and convexity give

\[
c_1d_k\le\Delta_k\le d_k.
\]

Expanding (7.2.1) at \(x_{k+1}=x_k-\alpha_kH_kg_k\) yields

\[
\boxed{
x_{k+1}^\top B_{k+1}x_{k+1}
=x_k^\top B_kx_k-\frac{f_k^2}{q_k}+\frac{h_k^2}{t_k}.
}
\tag{7.5.1}
\]

The explicit terms involving \(\alpha_k\) cancel because
\(g_k^\top x_k=f_k\):

\[
-2\alpha_k f_k+\alpha_k^2q_k
-\frac{(f_k-\alpha_kq_k)^2}{q_k}
=-\frac{f_k^2}{q_k}.
\]

Also,

\[
\frac{h_k^2}{t_k}\le h_k=d_k-\Delta_k
\le\frac{1-c_1}{c_1}\Delta_k.
\]

Consequently, with \(\gamma=(1-c_1)/c_1\),

\[
\Phi_k=x_k^\top B_kx_k+\gamma f_k,
\qquad
\boxed{\Phi_{k+1}\le\Phi_k-\frac{f_k^2}{q_k}.}
\tag{7.5.2}
\]

This potential inequality was established in an earlier source note. The new
ingredient is Theorem~\ref{M:average}, which replaces its former bounded-step
reciprocal-sum argument.

% SOURCE_END H.span.10


% SOURCE_BEGIN H.span.11
\subsection{Every finite value limit is
zero}\label{H:every-finite-value-limit-is-zero}

Armijo gives strictly decreasing values. Either they tend to \(-\infty\)
or they have a finite limit \(L\).

In the latter case,

\[
\sum_k d_k\le\frac{f_0-L}{c_1}<\infty,
\]

so \(d_k\to0\) and Theorem~\ref{M:average} gives \(\liminf q_k=0\). On the other hand,
\(\Phi_k\ge\gamma L\), and telescoping (7.5.2) gives

\[
\sum_k\frac{f_k^2}{q_k}<\infty.
\]

If \(L\ne0\), this would imply \(\sum_k1/q_k<\infty\) and therefore
\(q_k\to\infty\), a contradiction.

There is also a direct pointwise argument. The potential bounds

\[
x_k^\top B_kx_k\le\Phi_0-\gamma L=:M<\infty.
\]

Metric Cauchy--Schwarz and Euler's identity imply

\[
f_k^2=(g_k^\top x_k)^2
\le(g_k^\top H_kg_k)(x_k^\top B_kx_k)
\le Mq_k.
\]

A finite nonzero limit would thus force \(q_k\) to be bounded away from
zero, again contradicting Theorem~\ref{M:average}. Hence \(L=0\), proving Theorem 7.A.
\(\square\)

The formal finite-value theorem uses the same energy identity in a shorter alternative arrangement. If \(L\ne0\), shift to a tail with \(|f_k|\ge|L|/2\). The actual homogeneous corrections are nonnegative and bounded by \(d_k\), so they are summable. Chapter~\ref{comp:M} gives divergent reciprocal sums, and the energy identity would bound \((|L|/2)^2\sum q_k^{-1}\) by the initial tail energy plus a finite correction sum, a contradiction. The main formal proof invokes the already available superquadratic theorem to obtain this divergence. The displayed weighted-potential budgets and pointwise metric Cauchy--Schwarz bound are also separately proved; their presence does not assert that the implementation of the principal theorem follows both arguments verbatim.

% SOURCE_END H.span.11


% SOURCE_BEGIN H.span.12
\subsection{Bounded-below and coercive
consequences}\label{H:bounded-below-and-coercive-consequences}


A positively homogeneous function bounded below cannot take a negative
value: otherwise \(f(tx)=tf(x)\to-\infty\). Thus its minimum value is
zero, attained at the origin, and Theorem 7.A gives unconditional
objective convergence.

If its sublevels are bounded, then \(f(x)>0\) for \(x\ne0\). Compactness
of the unit sphere gives \(f(x)\ge\mu\|x\|\) for some \(\mu>0\), so
\(f_k\to0\) forces \(x_k\to0\).

These consequences now hold for all finite convex sublinear objectives,
not only finite polyhedral maxima. No rate or finite-path-length
estimate is asserted.


\setcounter{section}{5}
\setcounter{theorem}{0}
\setcounter{lemma}{0}
\setcounter{corollary}{0}
\setcounter{proposition}{0}
\setcounter{remark}{0}
\setcounter{equation}{0}
% SOURCE_END H.span.12


% SOURCE_BEGIN H.span.13
\section{Affine-hull obstruction and a many-piece escape
theorem}\label{H:affine-hull-obstruction-and-a-many-piece-escape-theorem}

\textbf{Lemma 7.3.} Suppose a full-BFGS run satisfies Armijo--weak-Wolfe,
has a finite gradient alphabet, and has objective values bounded below
along the run. Let \(G\) be the recurrent gradients. Then
\(0\in\operatorname{aff}G\).

\textbf{Proof.} Finite objective decrease gives \(d_k\to0\). Bounded
gradients and Theorem~\ref{M:average} give \(\liminf q_k=0\), hence
\(\sum_k1/q_k=\infty\).

Discard a finite prefix so that all remaining gradients belong to \(G\).
If \(0\notin\operatorname{aff}G\), a vector \(v\) exists with

\[
v^\top g=1\quad(g\in G).
\]

In particular \(v^\top y_k=0\). Equation (7.2.1) then implies

\[
v^\top B_{k+1}v=v^\top B_kv-\frac1{q_k}.
\]

Positive definiteness forces the reciprocal sum to be finite, a
contradiction. \(\square\)

This lemma does not need homogeneity, or even convexity, beyond what may
be needed to obtain the specified run. It is a necessary condition for
any finite-gradient, finite-value-limit trajectory under these BFGS and
line-search assumptions.

For max-linear objectives, if \(0\notin\operatorname{aff}\{a_i\}\), no
recurrent subset can have an affine hull containing zero. A finite value
limit is impossible; monotonicity therefore forces \(f_k\to-\infty\).
This proves Theorem 7.B.

% SOURCE_END H.span.13


% SOURCE_BEGIN H.span.14
\subsection{Connection to the two-piece
paper}\label{H:connection-to-the-two-piece-paper}


The model \(|u|+t\) has gradients \((1,1)\) and \((-1,1)\), both with
last component one. The arbitrary-many-piece model

\[
f(u,t)=t+\max_i b_i^\top u
\]

has precisely the same common-component property. The new theorem proves
that an infinite accepted-iterate sequence cannot stagnate on any finite
value for this entire family.

Xie--W\"achter~\cite{XW} proves the stronger finite-termination conclusion
for its two-piece model: a search direction exposing unboundedness is
reached after finitely many iterations. The argument here does
\textbf{not} establish finite termination for the general many-piece
family. It does show that any infinite admissible continuation must
decrease the values without bound.

If \(r=\dim\operatorname{span}G\), a finite-value-limit run needs at
least \(r+1\) distinct recurrent gradients. Indeed,
\(0\in\operatorname{aff}G\) implies that the gradient differences span
the same \(r\)-dimensional space; \(|G|-1\) differences must span it.
The three-gradient, rank-two stagnation example of Chapter~\ref{comp:D} attains this
lower bound.


\setcounter{section}{6}
\setcounter{theorem}{0}
\setcounter{lemma}{0}
\setcounter{corollary}{0}
\setcounter{proposition}{0}
\setcounter{remark}{0}
\setcounter{equation}{0}
% SOURCE_END H.span.14


% SOURCE_BEGIN H.span.15
\section{Iterate convergence in effective dimension at most
two}\label{H:iterate-convergence-in-effective-dimension-at-most-two}

% SOURCE_END H.span.15


% SOURCE_BEGIN H.span.16
\subsection{Exact invariant subspace and
lifting}\label{H:exact-invariant-subspace-and-lifting}

Suppose that, after some index \(K\), all gradients lie in a subspace
\(S\). Let \(U\) have orthonormal columns spanning \(S\), and define

\[
\widehat x_k=U^\top x_k,\qquad
\widehat H_k=U^\top H_kU,\qquad
\widehat g_k=U^\top g_k.
\]

Projection of the inverse-BFGS formula gives exactly the same
inverse-BFGS update for the hatted variables and projected steps. The
scalars \(q_k,d_k,t_k\) and the line-search inequalities are unchanged.

Moreover, all subsequent steps belong to the \textbf{fixed} subspace
\(H_KS\). To see this, initially \(H_KS\) is the image of \(S\); if
\(H_kS\subset H_KS\), then \(s_k\in H_KS\), and the inverse update maps
\(S\) back into \(H_KS\), since \(g_k,y_k\in S\). Induction proves the
claim.

The restriction of \(U^\top\) to \(H_KS\) is invertible because
\(U^\top H_KU\succ0\). Therefore

\[
\boxed{
x_k=x_K+L_K(\widehat x_k-\widehat x_K),\qquad
L_K=H_KU(U^\top H_KU)^{-1}.
}
\tag{7.7.1}
\]

In particular, convergence of the projected iterates implies convergence
of the original iterates. If \(\widehat x_k\to0\), the limit is

\[
\boxed{x_*=x_K-H_KU(U^\top H_KU)^{-1}U^\top x_K.}
\tag{7.7.2}
\]

The projected Hessian is \((U^\top H_kU)^{-1}\), not \(U^\top B_kU\).
This distinction is necessary for a valid reduction.

% SOURCE_END H.span.16


% SOURCE_BEGIN H.span.17
\subsection{A two-gradient endpoint
observation}\label{H:a-two-gradient-endpoint-observation}

Assume \(f_k\to0\). Then \(d_k\to0\), and every endpoint satisfies

\[
g_{k+1}^\top x_{k+1}=f_{k+1}\longrightarrow0,
\qquad
g_k^\top x_{k+1}=f_k-d_k\longrightarrow0.
\tag{7.7.3}
\]

In two dimensions, if \(g_k\) and \(g_{k+1}\) are linearly independent,
these two equations determine \(x_{k+1}\). There are only finitely many
distinct pairs, so their inverse norms are uniformly bounded. Thus the
endpoints of all independent-gradient transitions converge to zero, even
if the full sequence was not known to be bounded.

An unchanged nonzero gradient cannot satisfy weak Wolfe on a strict
descent step. Distinct gradients that are positive multiples of one
another also cannot both be active at positive-value iterates: among
positive multiples of a slope, only the largest multiple can attain a
positive maximum. Consequently, any dependent transition on the
positive-value branch is between opposite slopes.

% SOURCE_END H.span.17


% SOURCE_BEGIN H.span.18
\subsection{The unbounded-below planar
case}\label{H:the-unbounded-below-planar-case}

If the global objective is unbounded below, its slope set cannot contain
an opposite pair \(a,-\kappa a\) with \(\kappa>0\): their maximum alone
is nonnegative everywhere. Every accepted gradient transition is
therefore linearly independent in the plane.

Equation (7.7.3) now gives \(x_k\to0\) directly. Combined with Theorem 7.A,
this proves Theorem 7.C. In the other branch, global Lipschitz continuity
gives \(|f(x)|\le M\|x\|\), so \(f_k\to-\infty\) implies
\(\|x_k\|\to\infty\).

For an unbounded-below max-linear objective there is no stationary
point: a stationary point of a convex function would be a global
minimizer, and no minimizer exists. Thus the zero limit is genuinely
nonstationary.

% SOURCE_END H.span.18


% SOURCE_BEGIN H.span.19
\subsection{The general planar finite-value
branch}\label{H:the-general-planar-finite-value-branch}

For completeness, opposing-gradient transitions can also be handled,
giving convergence for \textbf{all} planar finite max-linear objectives,
including noncompact bounded-below ones.

Call an endpoint of an independent-gradient transition an \emph{anchor}.
By (7.7.3), anchors converge to zero.

Consider consecutive anchors \(x_a,x_b\). If there is no intermediate
iterate, there is nothing to prove. Otherwise, all gradients from
\(g_a\) through \(g_{b-1}\) lie on a single gradient line; the
intermediate dependent transitions alternate between opposite slopes. By
the invariant-subspace argument, all steps from \(x_a\) up to and
including the step arriving at \(x_b\) lie on a fixed primal line.
Therefore all points in this block are collinear with its endpoints.

Write its initial gradient as \(g\), and the next opposite gradient as
\(-\kappa g\), with \(\kappa>0\). The first dependent transition has

\[
d_a=f_a+f_{a+1}/\kappa.
\]

Armijo implies

\[
f_{a+1}\le\frac{1-c_1}{1+c_1/\kappa}f_a
\le(1-c_1)f_a,
\]

so monotonicity gives \(f_b\le(1-c_1)f_a\). Every accepted point in the
block obeys

\[
-f_a/\kappa\le g^\top x\le f_a,
\]

because \(g\) and \(-\kappa g\) are both slopes of the objective and
\(f(x)\le f_a\). The endpoint difference satisfies

\[
g^\top(x_b-x_a)\le-c_1 f_a.
\]

Writing \(x=x_a+\theta(x_b-x_a)\) consequently gives

\[
0\le\theta\le\frac{1+1/\kappa}{c_1}.
\]

Only finitely many slope pairs occur, so this bound is uniform across
blocks. Since both anchor endpoints tend to zero, every intermediate
point does too. Arbitrarily long blocks cause no difficulty.

If there are only finitely many anchors, the tail gradients lie on a
fixed line, and all tail iterates lie on one fixed primal line. The
slope coordinate tends to zero because \(f_k\to0\) and only finitely
many nonzero slope multipliers occur. Positive definiteness makes that
slope coordinate nonconstant along the primal line, so the entire tail
converges to its unique intersection with the zero-slope hyperplane.

This proves planar iterate convergence in every finite-value-limit case.
Reducing to the span of the recurrent gradients and using (7.7.1) proves
the same conclusion whenever that span has dimension at most two in any
ambient dimension. The limit may be nonzero in the original coordinates.


\setcounter{section}{7}
\setcounter{theorem}{0}
\setcounter{lemma}{0}
\setcounter{corollary}{0}
\setcounter{proposition}{0}
\setcounter{remark}{0}
\setcounter{equation}{0}
% SOURCE_END H.span.19


% SOURCE_BEGIN H.span.20
\section{A noncompact convergence theorem in arbitrary
dimension}\label{H:a-noncompact-convergence-theorem-in-arbitrary-dimension}


One source of noncompactness is a subspace of directions on which the
function is constant. When this is the only source of noncompactness,
the case is fully tractable in all dimensions.

Let

\[
S=\operatorname{span}\{a_1,\ldots,a_m\},\qquad
P=\operatorname{conv}\{a_1,\ldots,a_m\},
\]

and suppose \(0\in\operatorname{relint}P\), with relative interior taken
in \(S\). Then \(f\) depends only on the projection to \(S\) and is
coercive there:

\[
f(x)\ge\mu\|U^\top x\|\quad\text{for some }\mu>0.
\]

Theorem 7.A implies \(U^\top x_k\to0\). The invariant subspace formula
with \(K=0\) gives

\[
\boxed{
x_k\longrightarrow
x_0-H_0U(U^\top H_0U)^{-1}U^\top x_0.
}
\tag{7.8.1}
\]

The limit is a minimizer, and in fact the projection of \(x_0\) onto
\(S^\perp\) in the \(H_0^{-1}\) metric. It is independent of the choices
of admissible accepted lengths. This includes all finite polyhedral
seminorms and noncoercive gauges whose minimizer set is a linear
subspace.

The same argument applies after a finite prefix if
\(0\in\operatorname{relint}\operatorname{conv}G\) for the recurrent
gradients. The projected tail objective is then coercive even if some
original, nonrecurrent slopes have disappeared. Formula (7.7.2) supplies
the limit.

For \(f\equiv0\), the algorithm is already stationary; there is no
nonzero-gradient infinite run. The projection formula above concerns a
nontrivial slope span.


\endgroup
% SOURCE_END H.span.20


\section*{Relation to earlier work}
The following comparisons are bibliographic scholarship, separate from the formal theorem assertions.

The Euclidean-norm consequence is already Guo--Lewis~\cite[Corollary~4.2]{GL18}, while Lewis--Overton~\cite[Section~5.1]{LO} analyzes the scalar absolute-value model with its stated line-search conventions. The theorem here treats the larger finite convex homogeneous class without assuming bounded iterates.

The precise prior result is Xie--W\"achter's 2017 preprint~\cite[Theorem~IV.1]{XW}. Its finite detection of unbounded descent in the two-piece model is stronger there than the infinite-run value conclusion proved here for the many-piece family.

The invariant-subspace argument extends the affine-invariance background~\cite[Theorem~II.1]{XW}. Exact prior publication of the resulting $H_0^{-1}$-metric selector for this noncompact homogeneous class remains unresolved, as for Chapter~\ref{comp:F}'s rectangular lifting formula. This uncertainty concerns attribution, not an additional hypothesis of the theorem.
